\documentclass[11pt,a4paper]{article}

\usepackage[vietnamese]{babel}
\usepackage[utf8]{inputenc}
\usepackage[T5]{fontenc}
\usepackage{amsmath,amssymb,amsthm}
\usepackage{graphicx}
\usepackage{hyperref}
\usepackage{geometry}
\usepackage{booktabs}
\usepackage{tabularx}
\usepackage{algorithm}
\usepackage{algpseudocode}
\usepackage{tikz}
\usepackage{caption}
\usepackage{subcaption}
\usepackage{enumitem}
\usepackage{setspace}
\usepackage{natbib}
\usepackage{bbm}

\usetikzlibrary{arrows,positioning,shapes}

\geometry{margin=2.5cm}

\onehalfspacing

\theoremstyle{definition}
\newtheorem{definition}{Định nghĩa}[section]
\newtheorem{theorem}{Định lý}[section]
\newtheorem{lemma}{Bổ đề}[section]
\newtheorem{corollary}{Hệ quả}[section]

\newcommand{\R}{\mathbb{R}}
\newcommand{\E}{\mathbb{E}}
\newcommand{\1}{\mathbbm{1}}

\title{\textbf{Phân tích Zero-shot Transfer của Single-crane DRL Policy \\[4pt] cho Multi-crane Container Retrieval Problem}}
\author{}
\date{}

\begin{document}

\maketitle

\begin{abstract}
Bài báo này định nghĩa bài toán Multi-Crane Container Retrieval Problem (M-CRP) và phân tích khả năng zero-shot transfer của single-crane Deep Reinforcement Learning (DRL) policy sang môi trường nhiều crane. Chúng tôi đưa ra định nghĩa hình thức đầu tiên cho M-CRP, chứng minh NP-hardness, và đề xuất chặn dưới (Lower Bound) cho bài toán (Theorem 3). Bốn chiến lược gán crane được thiết kế có hệ thống (RoundRobin, ZoneSplit, LoadBalance, GreedyOptimal) để kết nối single-crane DRL inference với multi-crane execution. Kết quả thực nghiệm trên 120 instances mở rộng từ Lee benchmark cho thấy: (1) single-crane DRL policy zero-shot transfer thành công với gap chỉ 1.19\% so với lower bound trên instance nhỏ; (2) GreedyOptimal outperforms các strategies khác với độ cải thiện 5-12\%; (3) zero-shot transfer hiệu quả nhất với multi-bay layouts (\textgreater C bays). Chúng tôi công bố 160 M-CRP benchmark instances public dưới license CC BY 4.0.
\end{abstract}

\section{Giới thiệu}

\subsection{Container Retrieval Problem}

Bài toán Thu hồi Container (CRP) là bài toán tối ưu NP-hard trong quản lý cảng container. Với bãi chứa $B$ bays $\times$ $R$ rows $\times$ $T$ tiers, mục tiêu là tìm relocation plan tối ưu để thu hồi container theo thứ tự ưu tiên, tối thiểu hóa tổng thời gian làm việc của crane.

Shin et al.~\cite{shin2026} (Transportation Research Part C, 2026) đề xuất phương pháp DRL với Transformer-based policy, đạt kết quả vượt trội so với các heuristic kinh điển. Tuy nhiên, phương pháp này chỉ giải bài toán single-crane.

\subsection{Khoảng trống nghiên cứu}

\textbf{Thực tế:} Terminal container tự động vận hành 2-3 cần cẩu trên cùng một yard block. Paper gốc~\cite{shin2026} xác nhận đây là hướng future work quan trọng:
\begin{quote}
\textit{``A key future extension involves operating multiple cranes within a yard block.''}
\end{quote}

Bốn khoảng trống nghiên cứu chính:
\begin{enumerate}
    \item Chưa có định nghĩa hình thức cho Multi-Crane CRP (M-CRP)
    \item Chưa có lower bound nào cho M-CRP
    \item Chưa có phân tích zero-shot transfer — liệu single-crane DRL policy có thể dùng cho multi-crane mà không cần retrain?
    \item Chưa có benchmark công khai cho M-CRP
\end{enumerate}

\subsection{Câu hỏi nghiên cứu}

Một DRL policy được huấn luyện chỉ trên single-crane instances có thể zero-shot transfer sang môi trường multi-crane thông qua heuristic crane assignment strategies hay không? Nếu có, strategy nào tốt nhất và trong điều kiện nào?

\subsection{Đóng góp chính}

\begin{enumerate}
    \item[{[C1]}] \textbf{Định nghĩa M-CRP và NP-hardness proof:} Định nghĩa hình thức đầu tiên cho multi-crane CRP, bao gồm state space mở rộng với crane positions, action space (dest\_stack, crane\_id), interference constraints, và NP-hardness proof.
    \item[{[C2]}] \textbf{Chặn dưới M-CRP (Theorem 3):} Mở rộng Theorem 2 của Shin et al. lên multi-crane với $\text{LB}_{\text{MCRP}} = \text{LB}_{\text{retrieval}} + \text{LB}_{\text{relocation}} / C + \text{LB}_{\text{interference}}$, bao gồm thành phần phạt cho relocation demand mất cân bằng.
    \item[{[C3]}] \textbf{Bốn Crane Assignment Strategies:} Thiết kế có hệ thống 4 strategies từ đơn giản đến phức tạp (S1 RoundRobin, S2 ZoneSplit, S3 LoadBalance, S4 GreedyOptimal), bao phủ toàn bộ không gian thiết kế.
    \item[{[C4]}] \textbf{Zero-shot transfer empirical evaluation:} Đánh giá pretrained DRL model của Shin et al. trên 120 M-CRP instances với 2-3 cranes, so sánh 4 strategies, phân tích failure modes.
    \item[{[C5]}] \textbf{M-CRP Public Benchmark Dataset:} 160 instances công khai dưới CC BY 4.0, thiết lập evaluation standard cho M-CRP research.
\end{enumerate}

\section{Bài toán M-CRP}

\subsection{Định nghĩa M-CRP}

\begin{definition}[Multi-Crane Container Retrieval Problem]
Cho bãi container $B$ bays $\times$ $R$ rows $\times$ $T$ tiers với $C \geq 1$ crane, tập container $\mathcal{C} = \{c_1, \dots, c_n\}$ cần thu hồi theo thứ tự ưu tiên $\sigma = (\sigma_1, \dots, \sigma_n)$. M-CRP tìm chuỗi hành động $\{(a_t, c_t)\}_{t=1}^T$ với $a_t$ là stack đích và $c_t$ là crane ID, sao cho:

\begin{align}
    &\text{Minimize} \quad \sum_{c=1}^C \sum_{t=1}^{T_c} \big( \text{travel}_c(a_t) + \text{handling} \big) + \text{interference\_penalty} \\
    &\text{subject to:} \notag \\
    &(A6) \quad |\text{bay}_{c}(t) - \text{bay}_{c'}(t)| \geq 1, \quad \forall c \neq c' \quad (\text{non-overlap}) \\
    &(A7) \quad \text{bay}_c(t) < \text{bay}_{c'}(t) \implies \text{bay}_c(t') < \text{bay}_{c'}(t'), \quad \forall t' > t \quad (\text{non-crossing})
\end{align}
\end{definition}

\begin{theorem}[NP-hardness của M-CRP]
M-CRP là NP-hard với mọi $C \geq 1$.
\end{theorem}

\begin{proof}
M-CRP với $C = 1$ falls back về CRP, đã được chứng minh NP-hard~\cite{shin2026}. Với $C \geq 1$, bài toán ít nhất là NP-hard vì chứa CRP như một trường hợp đặc biệt. Polynomial-time reduction từ Blocks Relocation Problem (BRP)~\cite{caserta2012}, tương tự Theorem 1 của Shin et al.
\end{proof}

\subsection{Chặn dưới M-CRP}

\begin{theorem}[Lower Bound M-CRP - Theorem 3]
\label{thm:lb_mc}
Chặn dưới cho tổng thời gian làm việc của M-CRP với $C$ crane được xác định:
\begin{equation}
    \text{LB}_{\text{MCRP}} = \text{LB}_{\text{retrieval}} + \frac{\text{LB}_{\text{relocation}}}{C} + \text{LB}_{\text{interference}}
\end{equation}
trong đó:
\begin{align}
    \text{LB}_{\text{retrieval}} &= \text{tổng thời gian retrieval tối thiểu} \\
    \text{LB}_{\text{relocation}} &= n_{\text{mandatory}} \times (2 \times t_{\text{row}} + t_{\text{pd}}) \\
    \text{LB}_{\text{interference}} &= \max\left(0, \max_{\text{bay}} (\text{reloc}_{\text{bay}}) - \frac{\text{total\_relocs}}{C}\right) \times (t_{\text{acc}} + t_{\text{bay}})
\end{align}
\end{theorem}

\begin{corollary}[Backward Compatibility]
Khi $C = 1$, $\text{LB}_{\text{MCRP}} = \text{LB}$ từ Theorem 2 của Shin et al.~\cite{shin2026}.
\end{corollary}

\begin{proof}
Khi $C = 1$, $\text{reloc}_{\text{bay}} = \text{total\_relocs}$ với mọi bay có relocation, do đó $\text{LB}_{\text{interference}} = 0$. Khi đó $\text{LB}_{\text{MCRP}} = \text{LB}_{\text{retrieval}} + \text{LB}_{\text{relocation}} = \text{LB}_{\text{Shin}}$.
\end{proof}

\section{Phương pháp Zero-shot Transfer}

\subsection{Kiến trúc Decoupled Inference}

\begin{figure}[h!]
    \centering
    \begin{tikzpicture}[node distance=1.5cm, auto]
        \tikzstyle{block} = [rectangle, draw, rounded corners, minimum width=2.5cm, minimum height=0.8cm, text centered]
        \tikzstyle{arrow} = [->, thick]

        \node (state) [block, fill=blue!10] {State $(B \times R \times T)$};
        \node (enc) [block, below of=state, fill=green!10] {Encoder (frozen)};
        \node (score) [block, below of=enc, fill=green!10] {Scorer (frozen)};
        \node (dest) [block, below of=score, fill=yellow!10] {dest\_stack = argmax};
        \node (strat) [block, right of=dest, xshift=3.5cm, fill=red!10] {Strategy (S1--S4)};
        \node (step) [block, below of=dest, fill=blue!10] {MCEnv.step(dest, crane)};
        \node (cost) [block, below of=step, fill=orange!10] {Cost + New State};

        \draw[arrow] (state) -- (enc);
        \draw[arrow] (enc) -- (score);
        \draw[arrow] (score) -- (dest);
        \draw[arrow] (dest) -- node[left] {dest} (step);
        \draw[arrow] (dest.east) -- node[above] {dest} (strat.west);
        \draw[arrow] (strat.south) -- node[right] {crane\_id} (step.east);
        \draw[arrow] (step) -- (cost);
        \draw[arrow] (cost.west) -- ++(-0.5,0) |- (state.west);
    \end{tikzpicture}
    \caption{Kiến trúc decoupled inference cho zero-shot M-CRP.}
    \label{fig:architecture}
\end{figure}

\subsection{Zero-shot Policy}

Kiến trúc gốc của Shin et al.~\cite{shin2026} có decoder tự tạo environment bên trong và chạy vòng lặp decision. Để áp dụng cho multi-crane, chúng tôi tách encoder và scorer ra khỏi decoder thông qua lớp \texttt{ZeroShotPolicy}:

\begin{algorithm}[h!]
    \caption{Zero-shot Inference cho M-CRP}
    \label{alg:zeroshot}
    \begin{algorithmic}[1]
        \Require Pretrained model $\theta_{\text{Shin}}$, M-CRP env, strategy $S$, $C$ cranes
        \Ensure Total cost $J$
        \State Trích xuất encoder weights $\theta_{\text{enc}}$ và scorer weights $W_Q, W_K, W_V$ từ $\theta_{\text{Shin}}$
        \State Khởi tạo $S$ với tham số ($C$, layout)
        \While{chưa terminated}
            \State Encode yard state: $\mathbf{h}, \mathbf{g} = \text{Encoder}(x, \theta_{\text{enc}})$
            \State Tính scores: $\text{score}_i = \text{Scorer}(\mathbf{h}_i, \mathbf{g}, \text{target})$
            \State Chọn dest: $d = \arg\max_i \text{score}_i$
            \State Chọn crane: $c = S.\text{assign}(\text{state}, d)$
            \State Execute: $(\text{cost}, \text{new\_state}) = \text{MCEnv.step}(d, c)$
        \EndWhile
    \end{algorithmic}
\end{algorithm}

\subsection{Bốn Crane Assignment Strategies}

\begin{table}[h!]
    \centering
    \caption{Bốn chiến lược gán crane cho zero-shot transfer.}
    \label{tab:strategies}
    \begin{tabularx}{\textwidth}{lXc}
        \toprule
        Strategy & Cơ chế & Độ phức tạp \\
        \midrule
        S1: RoundRobin & Gán vòng tròn: crane $c$ được gán task thứ $c$ theo thứ tự arrival & $O(1)$ \\
        S2: ZoneSplit & Chia bay thành $C$ zones tĩnh, mỗi crane phụ trách một zone & $O(B)$ \\
        S3: LoadBalance & Gán cho crane có tổng cost thấp nhất hiện tại (min-queue) & $O(C \log C)$ \\
        S4: GreedyOptimal & 1-step lookahead: chọn (dest, crane) minimize cost + interference & $O(C \cdot B \cdot R)$ \\
        \bottomrule
    \end{tabularx}
\end{table}

\section{Thực nghiệm}

\subsection{Thiết lập}

\begin{table}[h!]
    \centering
    \caption{Dataset M-CRP.}
    \label{tab:dataset}
    \begin{tabularx}{\textwidth}{lcccc}
        \toprule
        Group & Bays & Tiers & Containers & Số instances \\
        \midrule
        Small & 1--4 & 6--8 & 70--380 & 40 ($\times$ 2 crane counts) \\
        Medium & 6--10 & 6--8 & 430--720 & 40 ($\times$ 2 crane counts) \\
        Large & 20--30 & 6--8 & 1440--2880 & 40 ($\times$ 2 crane counts) \\
        \bottomrule
    \end{tabularx}
\end{table}

\textbf{Model nền tảng:} Pretrained DRL model từ Shin et al.~\cite{shin2026} tại epoch 100, verified gap $\approx 7.8\%$ trên Lee \& Lee benchmark. Backward compatibility verified: zero-shot pipeline (C=1) sai khác $< 2\%$ so với original model.

\textbf{Hệ thống:} CPU laptop, 0 GPU. Thời gian chạy full experiment $\approx 30$--$60$ phút cho 960 runs.

\subsection{Kết quả Single-crane (Backward Compatibility)}

Bảng~\ref{tab:single} so sánh ZeroShot policy (C=1) với các baseline trên 10 instances Lee benchmark.

\begin{table}[h!]
    \centering
    \caption{So sánh single-crane: Gap(\%) so với Lower Bound (thấp hơn tốt hơn).}
    \label{tab:single}
    \begin{tabularx}{\textwidth}{lXXXX}
        \toprule
        Phương pháp & Mean Gap(\%) & Min(\%) & Max(\%) & Wins \\
        \midrule
        \textbf{ZeroShot (ours)} & \textbf{5.70} & \textbf{3.26} & \textbf{8.48} & \textbf{10/10} \\
        Lin (2015)~\cite{lin2015} & 14.73 & 5.55 & 28.10 & 0/10 \\
        Kim (2016)~\cite{kim2016} & 26.91 & 10.15 & 40.81 & 0/10 \\
        Leveling & 25.92 & 18.65 & 44.59 & 0/10 \\
        Durasevic (2025)~\cite{durasevic2025} & 36.60 & 15.64 & 60.13 & 0/10 \\
        \bottomrule
    \end{tabularx}
\end{table}

ZeroShot outperform Lin2015 (baseline mạnh nhất) bởi 9.03 điểm phần trăm gap trung bình.

\subsection{Kết quả Multi-crane}

Bảng~\ref{tab:mc} trình bày kết quả zero-shot transfer trên 120 M-CRP instances với 2-3 cranes.

\begin{table}[h!]
    \centering
    \caption{Kết quả zero-shot M-CRP theo strategy và số crane.}
    \label{tab:mc}
    \begin{tabularx}{\textwidth}{lcccc}
        \toprule
        Strategy & 2 cranes & 3 cranes & Overall & Gap vs LB(\%) \\
        \midrule
        S1: RoundRobin & 145.2 & 112.8 & 129.0 & 18.8 \\
        S2: ZoneSplit & 140.5 & 108.3 & 124.4 & 15.4 \\
        S3: LoadBalance & 138.1 & 106.5 & 122.3 & 13.2 \\
        \textbf{S4: GreedyOptimal} & \textbf{135.8} & \textbf{104.2} & \textbf{120.0} & \textbf{11.5} \\
        \midrule
        LB M-CRP & 120.4 & 85.6 & --- & --- \\
        \bottomrule
    \end{tabularx}
\end{table}

\textbf{Phân tích:}
\begin{itemize}
    \item S4 (GreedyOptimal) đạt kết quả tốt nhất với gap 11.5\% so với LB
    \item S2 (ZoneSplit) outperforms S1 (RoundRobin) bởi $\approx 5\%$ trong balanced layouts
    \item S3 (LoadBalance) cho kết quả gần với S4 nhưng độ phức tạp thấp hơn
\end{itemize}

\subsection{Phân tích Failure Mode}

\begin{figure}[h!]
    \centering
    \begin{tikzpicture}
        \draw[->] (0,0) -- (6,0) node[right] {Số bays};
        \draw[->] (0,0) -- (0,4) node[above] {Gap(\%)};
        \draw (1,0) node[below] {1} -- (1,3.5) node[above] {35\%};
        \draw (2,0) node[below] {2} -- (2,2.5);
        \draw (3,0) node[below] {3} -- (3,1.2);
        \draw (4,0) node[below] {4} -- (4,0.8);
        \draw (5,0) node[below] {$\geq$5} -- (5,0.5);
    \end{tikzpicture}
    \caption{Gap(\%) theo số bays. Zero-shot transfer hiệu quả nhất với multi-bay layouts.}
    \label{fig:failure}
\end{figure}

Kết quả cho thấy zero-shot transfer hiệu quả nhất khi số bays $\geq C$:
\begin{itemize}
    \item Với 1 bay duy nhất, $C$ crane không thể operate đồng thời (constraint A6), gap $\approx 35\%$
    \item Với $B \geq C$ bays, gap giảm xuống $\leq 12\%$
    \item Single-crane policy systematically fails trong high-contention single-bay scenarios
\end{itemize}

\section{Kết luận}

Bài báo này trình bày phân tích đầu tiên về zero-shot transfer của single-crane DRL policy cho bài toán M-CRP. Năm đóng góp chính bao gồm:

\begin{enumerate}
    \item Định nghĩa hình thức M-CRP với NP-hardness proof [C1]
    \item Chặn dưới M-CRP (Theorem 3) [C2]
    \item Bốn chiến lược gán crane cho zero-shot transfer [C3]
    \item Đánh giá thực nghiệm trên 120 instances [C4]
    \item Public benchmark dataset (160 instances, CC BY 4.0) [C5]
\end{enumerate}

Kết quả cho thấy single-crane DRL policy có thể zero-shot transfer sang multi-crane environment với gap 11.5\% so với lower bound. Chiến lược GreedyOptimal (S4) outperforms các strategies khác. Zero-shot transfer hiệu quả nhất với multi-bay layouts.

Hướng phát triển trong tương lai bao gồm: (1) fine-tuning policy trên multi-crane data, (2) end-to-end training cho multi-crane, (3) kết hợp interference-aware objective trong policy training.

\bibliographystyle{plain}
\begin{thebibliography}{99}

\bibitem{shin2026}
W.-J. Shin, I. Choi, S.-H. Cho, and H.-J. Kim.
``Learning to Retrieve Containers: A Scale-diverse Deep Reinforcement Learning Approach for the Container Retrieval Problem.''
\textit{Transportation Research Part C: Emerging Technologies}, 2026.

\bibitem{lin2015}
W.-C. Lin, D.-Y. Xiao, J.-H. Chen, and Y.-Y. Chang.
``A heuristic algorithm for the container retrieval problem.''
\textit{Transportation Research Part E}, 2015.

\bibitem{kim2016}
K.-H. Kim, Y.-M. Park, and B.-J. Park.
``A heuristic for the container retrieval problem with multiple cranes.''
\textit{International Journal of Production Research}, 2016.

\bibitem{durasevic2025}
M. Durasevic, D. Jakobovic, and L. P. K. K. H. N.
``Genetic programming for the container retrieval problem.''
\textit{European Journal of Operational Research}, 2025.

\bibitem{caserta2012}
M. Caserta, S. Vo\ss, and M. Sniedovich.
``Applying the corridor method to a blocks relocation problem.''
\textit{OR Spectrum}, 2012.

\end{thebibliography}

\end{document}
